
\renewcommand{\arraystretch}{1.5} % Default value: 1
%%%%%%%%%%%%%%
% Introduction
\section{Introduction}
The Mach-Zehnder interferometer is a simple but powerful tool used to show how light waves interfere with each other. It works by splitting a single beam of light into two separate paths using a "beamsplitter". These two beams travel along different routes before being brought back together by a second beamsplitter. Depending on what happens to the light on those two paths, the beams will either add up (constructive interference) or cancel each other out (destructive interference) when they reach the detectors.
%%%%%%%%%%%%%%%%%%%%%%%%%%%%%%%%%%%%%
\section{Objective}

%%%%%%%%%%%%%%%%%%%%%%%%%%%%%%%%%%%%%
\section{Equipment Required}
The following equipments are required to perform this experiment:
\begin{itemize}
    \item $650 \text{ nm}$ wavelength Laser source,
    \item Two 50-50 Beam splitters,
    \item Two reflectors (mirrors),
    \item Two lenses for focusing,
    \item Two white screens for marking and measuring the shift of the interference fringes,
    \item The optical bench, optical mounts and holders,
    \item A lens for producing circular fringes,
    \item A vacuum cell, and a vacuum gauge with a hand pump.
\end{itemize}
%%%%%%%%%%%%%%%%%%%%%%%%%%%%%%%%%%%%%
\section{Theory}
%%%%%%%%%%%%%%%%%%%%%%%%%%%%%%%%%%%%%
\section{Procedure}

%%%%%%%%%%%%%%%%%%%%%%%%%%%%%%%%%%%%%
\section{Data analysis and calculation} 
%%%%%%%%%%%%%%%%%%%%%%%%%%%%%%%%%%%%%
\section{Conclusion}
%%%%%%%%%%%%%%%%%%%%%%%%%%%%%%%%%%%%%